A linear variety is isomorphic to $\mathbb{P}^r$ by a projective linear change of coordinates, so its Hilbert polynomial is $\binom{z+r}{r}$ and its degree is $1$.

Conversely, degree is the sum of the degrees of the top-dimensional irreducible components. Since $Y$ is pure-dimensional and has degree $1$, it is irreducible. If $r=0$, it is a point. If $r=1$ and $Y$ is not a line, choose distinct $P,Q\in Y$ and a hyperplane through them which does not contain $Y$. Theorem 7.7 makes its intersection with $Y$ have degree $1$, whereas it contains two distinct points, a contradiction.

Proceed by induction on $r$. If $Y=\mathbb{P}^n$, there is nothing to prove. Otherwise, for any distinct $P,Q\in Y$, choose a hyperplane through $P,Q$ not containing $Y$; such a hyperplane exists unless $Y$ is contained in the line $PQ$. Its intersection with $Y$ has pure dimension $r-1$ and degree $1$ by Theorem 7.7, since all intersection multiplicities are positive. By induction the intersection is linear, so it contains the line $PQ$. Thus $Y$ contains every line joining two of its points. Its affine cone is consequently closed under addition and scalar multiplication, hence is a vector subspace. Therefore $Y$ is linear.
